\SourceLeaf{385}%
\textsc{von} \textsc{Zach} dort 1809 durchgeführten Triangulierung deutlicher erkennen lässt.

Im folgenden Abschnitt findet man verschiedene interessante Beobachtungen und Untersuchungen zum Zustand der Atmosphäre. Es sind dies gleichzeitige meteorologische Beobachtungen im Hospiz auf dem Mont Cenis und in Mailand, eine barometrische Höhenbestimmung des ersteren Orts und des Mont Colombier, eine trigonometrische Höhenbestimmung des Mont-blanc (4802,7 Meter) und des Monte Rosa (4619,6 Meter) und schließlich Untersuchungen zur terrestrischen Refraktion. Diese wird aus den Höhenwinkeln dreier Berge, Rocca Melone, Monte Viso und Monte Rosa, bestimmt, die an drei Punkten (Mailand, Turin, Mondovi) beobachtet wurden. Dabei werden die absoluten Höhen der drei Standpunkte als bekannt vorausgesetzt und die Höhen der beobachteten Punkte eliminiert. Dieses Verfahren kann wenig Sicherheit bieten und hat, wie eine genauere Prüfung zeigt, auch wenig Übereinstimmung ergeben. Wir möchten daher dem Endergebnis für das Verhältnis der Erdkrümmung zur \textit{gesamten} Refraktion (5,28 zu 1) wenig Gewicht beimessen. Die Rechnung aus sechs Paaren gegenseitiger Zenitdistanzen zwischen Hauptdreieckspunkten ergibt dieses Verhältnis im Mittel viel kleiner, nämlich 1 zu 0,1235, in sehr enger Übereinstimmung mit den Ergebnissen der hannoverschen und livländischen Gradmessung.

Der zehnte Abschnitt befasst sich mit der oft behandelten Aufgabe, aus der Breite des einen Endpunkts einer gegebenen Dreiecksseite und ihrem Azimut an diesem Endpunkt dieselben Größen für den anderen Endpunkt sowie den Längenunterschied auf dem elliptischen Sphäroid zu finden. Die Herleitung enthält nur eine Umformung der Formeln von \textsc{Legendre}, um die wahre Breite anstelle der sogenannten reduzierten Breite einzuführen. Die Formeln sind bis einschließlich der Größen dritter Ordnung genau, sofern man die Abplattung und die Dreiecksseite (den Erdradius als Einheit genommen) als Größen erster Ordnung betrachtet. Bei der Anwendung auf die vorliegenden Messungen hat man die Größen dritter Ordnung weggelassen, weil das für die Praxis genau genug sei. Das kann man allerdings nur zugestehen, sofern man die Ergebnisse lediglich zum Vergleich mit den astronomischen Bestimmungen verwenden will. Dazu ist es in der Tat unnötig, jene mit viel größerer Genauigkeit zu berechnen, als diese zulassen. Geht man aber von einem anderen Gesichtspunkt aus und berechnet die geodätischen Ergebnisse so genau, wie es
